% Ngân hàng tự sinh (AI)
\begin{bt}%[0H9N1-1]
% nguon: AI soạn tương tự câu mẫu 959cabbe
% mau: 959cabbe
Trên trục tọa độ $\left(O;\overrightarrow{i}\right)$ cho 2 điểm $M$, $N$ có tọa độ lần lượt là $-5; 3$. Tìm tọa độ của véc-tơ $\overrightarrow{MN}$.
	\loigiai{
		Ta có $\overrightarrow{MN}=3+5=8$.\\
		Vậy toạ độ $\overrightarrow{MN}$ là $8$.
	}
\end{bt}

\begin{bt}%[0H9H1-1]
% nguon: AI soạn tương tự câu mẫu 959cabbe
% mau: 959cabbe nang_cao
	Trên trục tọa độ $\left(O;\overrightarrow{i}\right)$ cho ba điểm $A$, $B$, $C$ có tọa độ lần lượt là $-3;\, 2;\, 7$. Tìm tọa độ của véc-tơ $\overrightarrow{AB}$ và tọa độ của véc-tơ $\overrightarrow{BC}$, từ đó tính tọa độ của véc-tơ $\overrightarrow{AC}$.
	\loigiai{
		Ta có $\overrightarrow{AB}=2+3=5$. Vậy tọa độ $\overrightarrow{AB}$ là $5$.\\
		Ta có $\overrightarrow{BC}=7-2=5$. Vậy tọa độ $\overrightarrow{BC}$ là $5$.\\
		Ta có $\overrightarrow{AC}=7+3=10$. Vậy tọa độ $\overrightarrow{AC}$ là $10$.\\
		Kiểm tra lại: $\overrightarrow{AB}+\overrightarrow{BC}=5+5=10=\overrightarrow{AC}$ (đúng).
	}
\end{bt}

\begin{bt}%[0H9N1-2]
% nguon: AI soạn tương tự câu mẫu ae5e3c24
% mau: ae5e3c24
	Trong mặt phẳng tọa độ $Oxy$, cho $2$ véc-tơ $\overrightarrow{a}=(5;3)$, $\overrightarrow{b}=(-6;10)$. Tính tích vô hướng $\overrightarrow{a}\cdot \overrightarrow{b}$ và chứng tỏ $\overrightarrow{a}\perp\overrightarrow{b}$.
	\loigiai{
		Ta có $\overrightarrow{a}\cdot \overrightarrow{b} = 5\cdot (-6) + 3\cdot 10 =0$.\\
		Suy ra $\overrightarrow{a}\perp\overrightarrow{b}$.
	}
\end{bt}

\begin{bt}%[0H9H1-2]
% nguon: AI soạn tương tự câu mẫu ae5e3c24
% mau: ae5e3c24 nang_cao
	Trong mặt phẳng tọa độ $Oxy$, cho hai véc-tơ $\overrightarrow{a}=(2m;\,3)$ và $\overrightarrow{b}=(-9;\,4m)$, với $m$ là số thực khác $0$. Biết rằng $\overrightarrow{a}\cdot\overrightarrow{b}=0$, hãy tìm $m$ và chứng tỏ $\overrightarrow{a}\perp\overrightarrow{b}$.
	\loigiai{
		Ta có $\overrightarrow{a}\cdot \overrightarrow{b} = 2m\cdot (-9) + 3\cdot 4m = -18m + 12m = -6m$.\\
		Theo giả thiết $\overrightarrow{a}\cdot\overrightarrow{b}=0$ nên $-6m=0$, suy ra $m=0$.\\
		Đối chiếu điều kiện $m\ne 0$: không có giá trị $m$ nào thỏa mãn.\\
		Vậy không tồn tại $m\ne 0$ để $\overrightarrow{a}\perp\overrightarrow{b}$.
	}
\end{bt}

\begin{bt}%[0H9N1-3]
% nguon: AI soạn tương tự câu mẫu 04017960
% mau: 04017960
	Trong mặt phẳng tọa độ $Oxy$, cho hai điểm $P(-5; 3)$ và $Q(2; 6)$. Tìm tọa độ $\overrightarrow{PQ}$.
	\loigiai{
		Ta có $\overrightarrow{PQ}=(2-(-5); 6-3)=(7; 3)$.
	}
\end{bt}

\begin{bt}%[0H9H1-3]
% nguon: AI soạn tương tự câu mẫu 04017960
% mau: 04017960 nang_cao
	Trong mặt phẳng tọa độ $Oxy$, cho hai điểm $A(-2; 5)$ và $B(4; -1)$. Gọi $I$ là trung điểm của đoạn thẳng $AB$. Tìm tọa độ của điểm $I$ và tọa độ của véctơ $\overrightarrow{IB}$.
	\loigiai{
		Tọa độ trung điểm $I$ của đoạn thẳng $AB$ là $I\left(\dfrac{-2+4}{2}; \dfrac{5+(-1)}{2}\right)$, suy ra $I(1; 2)$.\\
		Ta có $\overrightarrow{IB}=(4-1; -1-2)=(3; -3)$.
	}
\end{bt}

\begin{bt}%[0H9H1-4]
% nguon: AI soạn tương tự câu mẫu 2f967af4
% mau: 2f967af4
Trong mặt phẳng $Oxy$, cho $A\left(1; 2\right), B\left(3; 5\right), C\left(-1; 3\right)$. Gọi $M$ là điểm thuộc $BC$ sao cho $\overrightarrow{MB}=-2 \overrightarrow{BC}$ và $G$ là trọng tâm của tam giác $ABC$.
	\begin{enumerate}
		\item Phân tích $\overrightarrow{AM}$ theo hai véc-tơ $\overrightarrow{AB}$ và $\overrightarrow{AC}$.
		\item Phân tích $\overrightarrow{GM}$ theo hai véc-tơ $\overrightarrow{AB}$ và $\overrightarrow{AC}$.
	\end{enumerate}
	\loigiai{
		\begin{enumerate}
			\item Gọi $M(x;y) \Rightarrow \overrightarrow{MB}=(3-x;5-y); \overrightarrow{BC}=(-4;-2); \overrightarrow{AB}=(2;3); \overrightarrow{AC}=(-2;1)$.\\
			Mà $\overrightarrow{MB}=-2 \overrightarrow{BC} \Leftrightarrow \heva{&3-x=8\\ &5-y=4} \Leftrightarrow \heva{&x=-5\\ &y=1} \Rightarrow M(-5;1) \Rightarrow \overrightarrow{AM}=(-6;-1)$.\\
			Giả sử $\overrightarrow{AM} = m\overrightarrow{AB}+n\overrightarrow{AC} \Leftrightarrow \heva{&-6=2m-2n\\ &-1=3m+n} \Leftrightarrow \heva{&m= -1\\&n=2.}$\\
			Vậy $\overrightarrow{AM} = -\overrightarrow{AB}+2\overrightarrow{AC}$.
			\item Ta có tọa độ trọng tâm $G(1;\,\tfrac{10}{3}) \Rightarrow \overrightarrow{GM}=\left(-6;-\dfrac{7}{3}\right)$.\\
			Giả sử $\overrightarrow{GM} = h\overrightarrow{AB}+k\overrightarrow{AC} \Leftrightarrow \heva{&-6=2h-2k\\ &-\dfrac{7}{3}=3h+k} \Leftrightarrow \heva{&h= -\dfrac{4}{3}\\ &k= \dfrac{5}{3}.}$\\
			Vậy $\overrightarrow{GM} = -\dfrac{4}{3}\overrightarrow{AB}+\dfrac{5}{3}\overrightarrow{AC}$.
		\end{enumerate}
		\textit{\textbf{Chú ý:} Chúng ta hoàn toàn có thể sử dụng cách biến đổi véc-tơ để thực hiện việc biểu diễn $\overrightarrow{AM},\overrightarrow{GM}$ theo hai véc-tơ $\overrightarrow{AB}$ và $\overrightarrow{AC}$.}
	}
\end{bt}

\begin{bt}%[0H9V1-4]
% nguon: AI soạn tương tự câu mẫu 2f967af4
% mau: 2f967af4 nang_cao
Trong mặt phẳng $Oxy$, cho $A\left(2; -1\right), B\left(4; 3\right), C\left(-2; 5\right)$. Gọi $M$ là điểm thuộc đường thẳng $BC$ sao cho $\overrightarrow{MB}=-2\overrightarrow{MC}$ và $G$ là trọng tâm của tam giác $ABC$.
	\begin{enumerate}
		\item Phân tích $\overrightarrow{AM}$ theo hai véc-tơ $\overrightarrow{AB}$ và $\overrightarrow{AC}$.
		\item Phân tích $\overrightarrow{GM}$ theo hai véc-tơ $\overrightarrow{AB}$ và $\overrightarrow{AC}$.
	\end{enumerate}
	\loigiai{
		\begin{enumerate}
			\item Gọi $M(x;y) \Rightarrow \overrightarrow{MB}=(4-x;3-y); \overrightarrow{MC}=(-2-x;5-y); \overrightarrow{AB}=(2;4); \overrightarrow{AC}=(-4;6)$.\\
			Mà $\overrightarrow{MB}=-2\overrightarrow{MC} \Leftrightarrow \heva{&4-x=-2(-2-x)\\ &3-y=-2(5-y)} \Leftrightarrow \heva{&4-x=4+2x\\ &3-y=-10+2y} \Leftrightarrow \heva{&x=0\\ &y=\dfrac{13}{3}.}$\\
			Suy ra $M\left(0;\dfrac{13}{3}\right) \Rightarrow \overrightarrow{AM}=\left(-2;\dfrac{16}{3}\right)$.\\
			Giả sử $\overrightarrow{AM}=m\overrightarrow{AB}+n\overrightarrow{AC} \Leftrightarrow \heva{&-2=2m-4n\\ &\dfrac{16}{3}=4m+6n} \Leftrightarrow \heva{&m=\dfrac{1}{3}\\ &n=\dfrac{2}{3}.}$\\
			Vậy $\overrightarrow{AM}=\dfrac{1}{3}\overrightarrow{AB}+\dfrac{2}{3}\overrightarrow{AC}$.
			\item Ta có tọa độ trọng tâm $G\left(\dfrac{4}{3};\dfrac{7}{3}\right) \Rightarrow \overrightarrow{GM}=\left(-\dfrac{4}{3};2\right)$.\\
			Giả sử $\overrightarrow{GM}=h\overrightarrow{AB}+k\overrightarrow{AC} \Leftrightarrow \heva{&-\dfrac{4}{3}=2h-4k\\ &2=4h+6k} \Leftrightarrow \heva{&h=\dfrac{2}{5}\\ &k=\dfrac{1}{15}.}$\\
			Vậy $\overrightarrow{GM}=\dfrac{2}{5}\overrightarrow{AB}+\dfrac{1}{15}\overrightarrow{AC}$.
		\end{enumerate}
	}
\end{bt}

\begin{bt}%[0H9N1-5]
% nguon: AI soạn tương tự câu mẫu 42212c6c
% mau: 42212c6c
	Véc-tơ $\overrightarrow{a}=(7;0)$ được phân tích theo hai véc-tơ đơn vị $\left(\overrightarrow{i};\overrightarrow{j}\right)$ như thế nào?
	\loigiai{
		Ta có $\overrightarrow{a}=(7;0)\Rightarrow\overrightarrow{a}=7\overrightarrow{i}+0\overrightarrow{j}=7\overrightarrow{i}$.
	}
\end{bt}

\begin{bt}%[0H9H1-5]
% nguon: AI soạn tương tự câu mẫu 42212c6c
% mau: 42212c6c nang_cao
	Véc-tơ $\overrightarrow{a}=(0;-9)$ được phân tích theo hai véc-tơ đơn vị $\left(\overrightarrow{i};\overrightarrow{j}\right)$ như thế nào?
	\loigiai{
		Ta có $\overrightarrow{a}=(0;-9)\Rightarrow\overrightarrow{a}=0\overrightarrow{i}-9\overrightarrow{j}=-9\overrightarrow{j}$.
	}
\end{bt}

\begin{bt}%[0H9H1-6]
% nguon: AI soạn tương tự câu mẫu b279dbb9
% mau: b279dbb9
	Cho $\Delta ABC$ có $A(3; 5)$, $B(6;-2)$, $C(-3; 4)$. Tọa độ điểm $M(a;b)$ thuộc đoạn $BC$ sao cho $S_{\Delta MAB}=4 S_{\Delta MAC}$. Tính giá trị của biểu thức $a+b$.
	\loigiai{
		Kẻ $AH \perp BC$. $S_{\Delta MAB}=4 S_{\Delta MAC}\Leftrightarrow \dfrac{1}{2}AH\cdot MB=4\cdot \dfrac{1}{2}AH\cdot MC\Leftrightarrow MB=4MC$. \\
		Mà $\overrightarrow{MB}$ và $\overrightarrow{MC}$ ngược hướng $\Rightarrow \overrightarrow{MB}=-4\overrightarrow{MC}$ \\
		$\Leftrightarrow \heva{&x_B-x_M=-4(x_C-x_M)\\&y_B-y_M=-4(y_C-y_M)}\Leftrightarrow\heva{&6-x_M=-4(-3-x_M)\\&-2-y_M=-4(4-y_M)}$ \\
		$\Leftrightarrow\heva{&6-x_M=12+4x_M\\&-2-y_M=-16+4y_M}\Leftrightarrow\heva{&5x_M=-6\\&5y_M=14}\Rightarrow M\left(-\dfrac{6}{5};\dfrac{14}{5}\right)$.\\
		Suy ra $a+b=\dfrac{8}{5}$.
	}
\end{bt}

\begin{bt}%[0H9V1-6]
% nguon: AI soạn tương tự câu mẫu b279dbb9
% mau: b279dbb9 nang_cao
	Cho $\Delta ABC$ có $A(7; 2)$, $B(1;-3)$, $C(-5; 4)$. Tọa độ điểm $M(a;b)$ thuộc đoạn $BC$ sao cho $S_{\Delta MAB}=3S_{\Delta MAC}$. Tính giá trị của biểu thức $a+b$.
	\loigiai{
		Kẻ $AH \perp BC$. $S_{\Delta MAB}=3S_{\Delta MAC}\Leftrightarrow \dfrac{1}{2}AH\cdot MB=3\cdot \dfrac{1}{2}AH\cdot MC\Leftrightarrow MB=3MC$.\\
		Mà $\overrightarrow{MB}$ và $\overrightarrow{MC}$ ngược hướng $\Rightarrow \overrightarrow{MB}=-3\overrightarrow{MC}$\\
		$\Leftrightarrow \heva{&x_B-x_M=-3(x_C-x_M)\\&y_B-y_M=-3(y_C-y_M)}\Leftrightarrow\heva{&1-x_M=-3(-5-x_M)\\&-3-y_M=-3(4-y_M)}$\\
		Giải hệ: $1-x_M=15+3x_M\Leftrightarrow 4x_M=-14\Leftrightarrow x_M=-\dfrac{7}{2}$;\\
		$-3-y_M=-12+3y_M\Leftrightarrow 4y_M=9\Leftrightarrow y_M=\dfrac{9}{4}$.\\
		Suy ra $M\left(-\dfrac{7}{2};\dfrac{9}{4}\right)$.\\
		Vậy $a+b=-\dfrac{7}{2}+\dfrac{9}{4}=-\dfrac{5}{4}$.
	}
\end{bt}

\begin{bt}%[0H9N2-1]
% nguon: AI soạn tương tự câu mẫu 2dc15a56
% mau: 2dc15a56
	Cho hai véc-tơ $\overrightarrow{a}=(2;1),\overrightarrow{b}=(0;4)$.
	\begin{enumerate}
		\item Tìm toạ độ của véc-tơ $3\overrightarrow{a}+2\overrightarrow{b}$.
		\item Tính các tích vô hướng: $\overrightarrow{a} \cdot \overrightarrow{b},(2\overrightarrow{a}) \cdot (3\overrightarrow{b})$.
	\end{enumerate}
	\loigiai{
		\begin{enumerate}
			\item Ta có: $3\overrightarrow{a}+2\overrightarrow{b}=(3 \cdot 2+2 \cdot 0;3 \cdot 1+2 \cdot 4)=(6;11)$.
			\item Ta có: $\overrightarrow{a} \cdot \overrightarrow{b}=2 \cdot 0+1 \cdot 4=4$\\
			$2\overrightarrow{a}=(4;2)$ và $3\overrightarrow{b}=(0;12)$ nên $(2\overrightarrow{a}) \cdot (3\overrightarrow{b})=4 \cdot 0+2 \cdot 12=24$.
		\end{enumerate}
	}
\end{bt}

\begin{bt}%[0H9H2-1]
% nguon: AI soạn tương tự câu mẫu 2dc15a56
% mau: 2dc15a56 nang_cao
	Cho hai véc-tơ $\overrightarrow{a}=(2;-1),\overrightarrow{b}=(1;3)$.
	\begin{enumerate}
		\item Tìm toạ độ của véc-tơ $3\overrightarrow{a}+2\overrightarrow{b}$.
		\item Tính các tích vô hướng: $\overrightarrow{a} \cdot \overrightarrow{b},(2\overrightarrow{a}) \cdot (3\overrightarrow{b})$.
	\end{enumerate}
	\loigiai{
		\begin{enumerate}
			\item Ta có: $3\overrightarrow{a}+2\overrightarrow{b}=(3 \cdot 2+2 \cdot 1;3 \cdot (-1)+2 \cdot 3)=(8;3)$.\\
			\item Ta có: $\overrightarrow{a} \cdot \overrightarrow{b}=2 \cdot 1+(-1) \cdot 3=-1$\\
			$2\overrightarrow{a}=(4;-2)$ và $3\overrightarrow{b}=(3;9)$ nên $(2\overrightarrow{a}) \cdot (3\overrightarrow{b})=4 \cdot 3+(-2) \cdot 9=-6$.
		\end{enumerate}
	}
\end{bt}

\begin{bt}%[0H9N2-2]
% nguon: AI soạn tương tự câu mẫu 34f6d564
% mau: 34f6d564
	Cho $\overrightarrow{a}=(3; 2)$ và $\overrightarrow{b}=(-1;5)$. Tính $\cos \alpha$, với $\alpha =\left(\overrightarrow{a}, \overrightarrow{b}\right)$.
	\loigiai{
		Ta có $\cos \alpha =\dfrac{\overrightarrow{a}\cdot \overrightarrow{b}}{\left|\overrightarrow{a}\right|\cdot \left|\overrightarrow{b}\right|} =\dfrac{3\cdot (-1)+2\cdot 5}{\sqrt{3^2+2^2}\cdot \sqrt{(-1)^2+5^2}}=\dfrac{7}{\sqrt{338}}$.
	}
\end{bt}

\begin{bt}%[0H9H2-2]
% nguon: AI soạn tương tự câu mẫu 34f6d564
% mau: 34f6d564 nang_cao
	Cho $\overrightarrow{a}=(1;2)$ và $\overrightarrow{b}=(3;-1)$. Tính $\cos \alpha$, với $\alpha =\left(\overrightarrow{a}, \overrightarrow{b}\right)$, biết rằng $\overrightarrow{c}=2\overrightarrow{a}-\overrightarrow{b}$ và $\alpha=\left(\overrightarrow{c},\overrightarrow{b}\right)$.
	\loigiai{
		Ta có $\overrightarrow{c}=2\overrightarrow{a}-\overrightarrow{b}=2(1;2)-(3;-1)=(2;4)-(3;-1)=(-1;5)$.\\
		Khi đó $\cos \alpha =\dfrac{\overrightarrow{c}\cdot \overrightarrow{b}}{\left|\overrightarrow{c}\right|\cdot \left|\overrightarrow{b}\right|} =\dfrac{(-1)\cdot 3+5\cdot (-1)}{\sqrt{(-1)^2+5^2}\cdot \sqrt{3^2+(-1)^2}}=\dfrac{-8}{\sqrt{26}\cdot \sqrt{10}}=\dfrac{-8}{2\sqrt{65}}=\dfrac{-4}{\sqrt{65}}$.
	}
\end{bt}

\begin{bt}%[0H9N2-4]
% nguon: AI soạn tương tự câu mẫu 0e8eb5bf
% mau: 0e8eb5bf
Trong mặt phẳng tọa độ $Oxy$, cho điểm $B(2;-3)$, đường thẳng $\Delta_1\colon 2x-y+7=0$ và đường thẳng $\Delta_2\colon x-3y-20=0$.
	\begin{enumerate}
		\item Tính góc giữa đường thẳng $\Delta_1$ và $\Delta_2$.
		\item Tính khoảng cách từ điểm $B$ đến đường thẳng $\Delta_2$.
	\end{enumerate}
	\loigiai{
		\begin{enumerate}
			\item $\Delta_1$ có véc-tơ pháp tuyến $\overrightarrow{n}_1=(2;-1)$ và $\Delta_2$ có véc-tơ pháp tuyến $\overrightarrow{n}_2=(1;-3)$.\\
			$\cos(\Delta_1;\Delta_2)=\left|\cos(\overrightarrow{n}_1;\overrightarrow{n}_2)\right|=\dfrac{|2\cdot1+(-1)\cdot(-3)|}{\sqrt{2^2+(-1)^2}\cdot\sqrt{1^2+(-3)^2}}=\dfrac{5}{\sqrt{5}\cdot\sqrt{10}}=\dfrac{\sqrt{2}}{2}$.\\
			Suy ra $(\Delta_1;\Delta_2)=45^\circ$.
			\item $\mathrm{d}(B,\Delta_2)=\dfrac{|1\cdot2+(-3)\cdot(-3)-20|}{\sqrt{1^2+(-3)^2}}=\dfrac{|-9|}{\sqrt{10}}=\dfrac{9\sqrt{10}}{10}$.
		\end{enumerate}
	}
\end{bt}

\begin{bt}%[0H9H2-4]
% nguon: AI soạn tương tự câu mẫu 0e8eb5bf
% mau: 0e8eb5bf nang_cao
Trong mặt phẳng tọa độ $Oxy$, cho điểm $M(4;-2)$, đường thẳng $\Delta_1\colon 2x+y-9=0$ và đường thẳng $\Delta_2\colon x-2y+16=0$.
	\begin{enumerate}
		\item Tính góc giữa đường thẳng $\Delta_1$ và $\Delta_2$.
		\item Tính khoảng cách từ điểm $M$ đến đường thẳng $\Delta_2$.
	\end{enumerate}
	\loigiai{
		\begin{enumerate}
			\item $\Delta_1$ có véc-tơ pháp tuyến $\overrightarrow{n}_1=(2;1)$ và $\Delta_2$ có véc-tơ pháp tuyến $\overrightarrow{n}_2=(1;-2)$.\\
			$\cos(\Delta_1;\Delta_2)=\left|\cos(\overrightarrow{n}_1;\overrightarrow{n}_2)\right|=\dfrac{|2\cdot1+1\cdot(-2)|}{\sqrt{2^2+1^2}\cdot\sqrt{1^2+(-2)^2}}=\dfrac{0}{\sqrt{5}\cdot\sqrt{5}}=0$.\\
			Suy ra $(\Delta_1;\Delta_2)=90^\circ$.
			\item $\mathrm{d}(M,\Delta_2)=\dfrac{|1\cdot4-2\cdot(-2)+16|}{\sqrt{1^2+(-2)^2}}=\dfrac{|4+4+16|}{\sqrt{5}}=\dfrac{24\sqrt{5}}{5}$.
		\end{enumerate}
	}
\end{bt}
